% from ega3-7-fr.tex
\oldpage[III]{222}
Posons
$\mathscr G_1=\operatorname{Im}(\mathscr F\to\mathscr G)
=\operatorname{Ker}(\mathscr G\to\mathscr H)$, qui est cohérent. Comme on a la
suite exacte
$0\to\mathscr G_1(n)\to\mathscr G(n)\to\mathscr H(n)$ et que le foncteur
$f_*$ est exact à gauche, la suite
$0\to f_*(\mathscr G_1(n))\to f_*(\mathscr G(n))\to f_*(\mathscr H(n))$ est
exacte pour tout $n$, et il suffit donc de montrer que, pour $n$ assez grand,
la suite $f_*(\mathscr F(n))\to f_*(\mathscr G_1(n))\to0$ est exacte. Autrement
dit, on peut se borner au cas où $\mathscr H=0$ et où le support de
$\mathscr G$ est propre sur $Y$, donc \emph{fermé dans $Z$} (II, 5.4.10) ;
$\mathscr G'=i_*(\mathscr G)$ est par suite un $\mathscr O_Z$-Module cohérent tel
que $\mathscr G'|X=\mathscr G$. On sait (I, 9.4.3) qu'il existe un
$\mathscr O_Z$-Module cohérent $\mathscr F'$ tel que
$\mathscr F'|X=\mathscr F$. En outre, comme $X$ est ouvert dans $Z$ et
$\operatorname{Supp}(\mathscr G)\subset X$ fermé dans $Z$, il est immédiat que
l'on définit un homomorphisme surjectif $u':\mathscr F'\to\mathscr G'$ de
faisceaux tel que $u'|X$ soit l'homomorphisme surjectif donné
$u:\mathscr F\to\mathscr G$, en prenant $u'|U=0$ pour tout ouvert $U$ de $Z$
ne rencontrant pas $\operatorname{Supp}(\mathscr G)$ et $u'|U=u|U$ pour tout
ouvert $U\subset X$. Cela étant, on voit comme au début de la démonstration de
(2.2.2) que l'on peut se borner au cas où $Y$ est affine, et il s'agit donc de
montrer que, pour $n$ assez grand, l'homomorphisme
$\Gamma(u):\Gamma(X,\mathscr F(n))\to\Gamma(X,\mathscr G(n))$ est surjectif. Or,
on a le diagramme commutatif
\[
\xymatrix@C=35mm@R=12mm{
\Gamma(Z,\mathscr F'(n)) \ar[r]^{\Gamma(u')} \ar[d]_v &
\Gamma(Z,\mathscr G'(n)) \ar[d]^w \\
\Gamma(X,\mathscr F(n)) \ar[r]_{\Gamma(u)} & \Gamma(X,\mathscr G(n))
}
\]
où $v$ et $w$ sont les homomorphismes de restriction. La définition de
$\mathscr G'$ montre en outre que $w$ est \emph{bijectif} ; par ailleurs, en
vertu de (2.2.3), $\Gamma(u')$ est \emph{surjectif} pour $n$ assez grand, donc
il en est de même de $\Gamma(u)$.

\noindent\textbf{$(\mathbf{Err}_{\mathrm{III}},\,27)$} Dans (III, 2.2.5),
après la ligne 7 du bas de la p. 102, ajouter :

(iii) Sous les hypothèses de (2.2.4) concernant $X$, $Y$, $f$ et $\mathscr L$,
il est immédiat que si $\mathscr H$ est un $\mathscr O_X$-Module cohérent dont
le support est \emph{propre sur $Y$}, un raisonnement analogue à celui de
(2.2.4) montre qu'il existe un entier $N$ tel que pour $n\geq N$, on ait
$\mathrm R^if_*(\mathscr H(n))=0$ pour tout $i>0$. On en conclut, par la suite
exacte de cohomologie, que si $u:\mathscr F\to\mathscr G$ est un homomorphisme
de $\mathscr O_X$-Modules cohérents tel que $\operatorname{Ker}(u)$ et
$\operatorname{Coker}(u)$ aient leurs supports \emph{propres sur $Y$}, alors il
existe $N$ tel que pour $n\geq N$, l'homomorphisme correspondant
$\mathrm R^if_*(\mathscr F(n))\to\mathrm R^if_*(\mathscr G(n))$ soit
\emph{bijectif} pour tout $i>0$.

\noindent\textbf{$(\mathbf{Err}_{\mathrm{III}},\,28)$} Dans (III, 4.4.9),
ligne 17 du bas de la p. 137, remplacer « \emph{Soient $Y$ un préschéma
intègre localement noethérien} » par « \emph{Soient $X$ et $Y$ deux
préschémas intègres localement noethériens} » ; ligne 16 du bas de la p.
137, remplacer « \emph{de type fini} » par « \emph{localement de type
fini} ». La démonstration est essentiellement inchangée,
$f^{-1}(y)$ étant localement de type fini sur $k(y)$, donc encore discret.

\noindent\textbf{$(\mathbf{Err}_{\mathrm{III}},\,29)$} dans (I, 9.3.4), ligne
19 du bas de la p. 173, remplacer « \emph{noethérien} » par
« \emph{quasi-compact} » ; lignes 18 et 19 du bas de la p. 173, remplacer
« \emph{cohérent} » par « \emph{quasi-cohérent de type fini} ». Dans la
démonstration, on se ramène au cas où

